% 復習4　文献を批判的に読む（記入例）
\session{4}{1}{90分}{文献を批判的に読む}{AI要約と対話}{}

\goals{
  \item 論文の構造（問い、方法、結果、考察、限界）を押さえて読める
  \item AI要約の使いどころと限界を説明できる。要約は原文の代わりにならない
  \item 主張と根拠の対応、データの範囲、別の解釈の可能性を確かめる批判的な読み方ができる
}

\afillin{読む文献（著者・年・題名）}{架空論文Ｙ「中小製造業における生成AIの試験導入と組織的課題」}

\begin{caution}{入力するときの注意}
論文の全文を入力する場合は、大学のAI利用規程と、その資料の利用条件を確認する。要旨や必要な部分だけを使えば十分なことが多い。
\end{caution}

\work[80mm]{ワーク1}{AI要約と原文を照らし合わせる}
\begin{promptbox}[構造をつかむ]
次の論文について、「問い」「方法」「主な結果」「著者が認めている限界」をそれぞれ2文以内でまとめてください。書かれていない項目は「記載なし」としてください。（論文の要旨や該当部分を貼る）
\end{promptbox}
\begin{wstabh}{colspec={Q[l,wd=17mm,m] X[l,m] X[l,m] X[l,m]},row{2-Z}={ht=17mm}}
  & AIの要約（要点） & 原文で確認した内容（頁） & 要約の誤り・抜け \\
  問い & \af{中小企業はなぜ生成AIを導入しないのか} & \af{生成AIを試験導入した中小製造業が、本格導入に進むかどうかを分ける要因は何か（p.2）} & \af{対象は「導入しない企業」ではなく「試験導入した企業」。要約が問いを広げている} \\
  方法 & \af{中小企業へのアンケート調査} & \af{試験導入した中小製造業12社への聞き取り調査（p.5）} & \af{アンケートではなく聞き取り。社数（12社）が抜けている} \\
  主な結果 & \af{費用が高いことが最大の障壁} & \af{費用より「使う業務を決められない」「社内で教える人がいない」が多く挙がった（p.9）} & \af{結果を取り違えている。費用は上位ではない} \\
  限界 & \af{記載なし} & \af{12社と少なく、地域も一つ。業種が製造業に限られる（p.14）} & \af{原文には限界の節がある。AIが見落とした} \\
\end{wstabh}

\work[70mm]{ワーク2}{最も強い反論に、原文の根拠で答える}
\begin{promptbox}[反論を引き出す]
この論文の主張「（主張）」に対して、最も強い反論を3つ挙げてください。それぞれについて、どのようなデータがあれば検証できるかも示してください。
\end{promptbox}
\begin{wstabh}{colspec={X[3,l,m] Q[c,wd=15mm,m] X[3,l,m]},row{2-Z}={ht=16mm}}
  AIが挙げた反論 & 答えられる\newline ○△× & 原文の根拠（頁）・答えられない理由 \\
  \af{12社では少なく、中小企業全体の傾向とは言えない} & \ans{×} & \af{原文も限界として認めている（p.14）。この反論には答えられない} \\
  \af{試験導入した企業は、もともとITに積極的な企業ではないか} & \ans{△} & \af{各社のIT投資の状況を記述している（表2、p.6）が、比べる相手（導入しない企業）がいない} \\
  \af{「使う業務を決められない」は、経営者の関心の低さの言い換えでは} & \ans{○} & \af{経営者自身が積極的な企業でも同じ課題が挙がっている（p.10）} \\
\end{wstabh}

\work{ワーク3}{未解決の要素}
\atwoboxes{論文が答えていないこと}{自分の問いとの隙間}{32mm}
{・試験導入をしていない企業は、なぜ導入しないのか\newline ・製造業以外ではどうか\newline ・「使う業務を決められない」を乗り越えた企業は何をしたのか}
{自分の問い（なぜ進まないか）は、導入しない企業まで含む広い問い。この論文に合わせて「試験導入した企業が、どの業務で使い、何を障壁と感じているか」に絞ると、調べやすくなる。}

\begin{nexttime}
  \item 読んだ文献について「主張」「根拠」「限界」を各100字程度でまとめる
  \item AIとの議論で見えた未解決の要素を1つ書き、自分の問いとどうつながるかを考える
\end{nexttime}
\abox[主張（100字程度）]{24mm}{生成AIを試験導入した中小製造業が本格導入に進めない主な理由は、費用よりも、使う業務を決められないことと、社内で使い方を教える人がいないことにある。導入には技術より組織の準備が重要である。}
\abox[根拠（100字程度）]{24mm}{試験導入した中小製造業12社に聞き取り調査を行い、本格導入に進んだ企業と進まなかった企業で、挙げられた課題を比べた。費用を最大の課題とした企業は少なく、業務の選定と教える人の不足が多かった。}
\abox[限界（100字程度）]{24mm}{対象が12社と少なく、地域も業種も限られるため、中小企業全体には一般化できない。試験導入していない企業は調べていないので、導入しない理由はわからない。}
\abox[未解決の要素と、自分の問いとのつながり]{24mm}{「使う業務を決められない」をどう乗り越えたかは書かれていない。自分の問いを試験導入した企業に絞り、どの業務で使い始めたかを聞くことで、この隙間を埋められる。}
\aimemoA{論文の構造の要約と、主張への反論の生成}
{「問い・方法・主な結果・限界を2文以内で。書かれていない項目は記載なしと」「最も強い反論を3つ」}
{要約は読む前の地図としてだけ使い、表の内容はすべて原文から書き直した。反論3つのうち2つは、自分の問いを絞り込む材料にした。}
{要約の各項目を原文の該当ページと照合した（誤りと抜けを表に記録）}

\begin{point}
  \item 記入例で使った論文は、この教材のための\textbf{架空の論文}。実際には、復習3で入手した文献を使う。
  \item AI要約には典型的なずれがある。問いを広げる、方法の条件（対象・人数）を落とす、結果を取り違える、限界を落とす。表の右の列で、これらを見つけられたかを確かめる。
  \item 原文の頁を書くと、原文を読んだ証拠になり、後で引用するときにも役立つ。
  \item 答えられない反論が出てくるのは悪いことではない。論文の範囲の外だと言えれば十分で、自分の問いを絞る手がかりになる。
\end{point}
